\documentclass[11pt]{article}

\usepackage[margin=1in]{geometry}
\usepackage{amsmath, amssymb}
\usepackage{algorithm}
\usepackage{algpseudocode}
\usepackage{booktabs}
\usepackage{hyperref}

\title{Kernel PCA}
\author{ML A-Z --- Part 11}
\date{}

\begin{document}
\maketitle

\section{Motivation}

Standard PCA (Part 10) finds directions of maximum variance by
eigendecomposing the covariance matrix $\Sigma = \frac{1}{n} X^\top X$.
Because $\Sigma$ is built from \emph{linear} combinations of the original
features, the components it finds are linear subspaces --- it can only
"see" variance along straight lines through the data. If the true
structure in the data is nonlinear (e.g.\ classes separated by a curved
boundary, or lying on a manifold like concentric circles), no linear
projection will reveal it, no matter how many components are kept.

\textbf{Kernel PCA} fixes this by first mapping the data into a much
higher-dimensional (possibly infinite-dimensional) feature space via a
nonlinear map $\phi$, where the structure becomes linear, and then running
ordinary PCA in that space. The obstacle is that $\phi$ can be expensive or
impossible to compute explicitly --- the \textbf{kernel trick} sidesteps
this entirely.

\section{The Kernel Trick}

\subsection{Idea}

Many linear algorithms, PCA included, only ever need the data through
inner products $\mathbf{x}_i^\top \mathbf{x}_j$, never through the raw
feature vectors themselves. The kernel trick observes that if an algorithm
can be written purely in terms of such inner products, then to run it in a
feature space $\phi(\mathbf{x})$ we don't need $\phi$ explicitly --- we
only need a function

\[
  k(\mathbf{x}_i, \mathbf{x}_j) = \phi(\mathbf{x}_i)^\top \phi(\mathbf{x}_j)
\]

that computes the inner product \emph{in that space directly from the
original vectors}. Replacing every $\mathbf{x}_i^\top \mathbf{x}_j$ with
$k(\mathbf{x}_i, \mathbf{x}_j)$ implicitly performs the algorithm in the
(possibly infinite-dimensional) space of $\phi$, at the cost of evaluating
$k$ instead of an explicit dot product --- often $O(p)$ work either way,
regardless of how large the target space is.

\subsection{Common Kernels}

\begin{center}
\begin{tabular}{lll}
  \toprule
  Kernel & Formula & Notes \\
  \midrule
  Linear & $k(\mathbf{x}, \mathbf{y}) = \mathbf{x}^\top \mathbf{y}$ & Recovers ordinary PCA \\
  Polynomial & $k(\mathbf{x}, \mathbf{y}) = (\gamma \, \mathbf{x}^\top \mathbf{y} + c)^d$ & Captures feature interactions up to degree $d$ \\
  RBF (Gaussian) & $k(\mathbf{x}, \mathbf{y}) = \exp(-\gamma \lVert \mathbf{x} - \mathbf{y} \rVert^2)$ & Infinite-dimensional feature space; most common default \\
  Sigmoid & $k(\mathbf{x}, \mathbf{y}) = \tanh(\gamma \, \mathbf{x}^\top \mathbf{y} + c)$ & Behaves like a neural-net activation \\
  \bottomrule
\end{tabular}
\end{center}

A function $k$ is a valid kernel --- i.e.\ actually corresponds to an inner
product in \emph{some} feature space --- iff it is symmetric and positive
semi-definite (\textbf{Mercer's condition}): for any finite set of points
$\mathbf{x}_1, \dots, \mathbf{x}_n$, the Gram matrix $K_{ij} = k(\mathbf{x}_i,
\mathbf{x}_j)$ must be PSD.

\section{From PCA to Kernel PCA}

\subsection{Ordinary PCA, Rewritten with Inner Products}

Assume the data is mean-centered in feature space. Ordinary PCA
diagonalizes $\Sigma = \frac{1}{n} X^\top X \in \mathbb{R}^{p \times p}$. It
turns out the same eigenvectors can be recovered from the
$n \times n$ Gram matrix $XX^\top$ instead: if $\mathbf{v}$ is a
(unit) eigenvector of $\Sigma$ with eigenvalue $\lambda$, then
$\mathbf{v} = \frac{1}{\sqrt{n\lambda}} X^\top \mathbf{a}$ for some
eigenvector $\mathbf{a}$ of $XX^\top$ with the same eigenvalue (up to a
factor of $n$), i.e.\ every principal direction is expressible as a linear
combination of the (centered) data points themselves. This is the key fact
that lets PCA be reformulated using only pairwise inner products
$X X^\top$, without ever referring to $\Sigma$ or the raw feature space
directly.

\subsection{Kernelizing}

Replace every inner product $\mathbf{x}_i^\top \mathbf{x}_j$ with
$k(\mathbf{x}_i, \mathbf{x}_j)$, giving the $n \times n$ kernel (Gram)
matrix $K$ with $K_{ij} = k(\mathbf{x}_i, \mathbf{x}_j) = \phi(\mathbf{x}_i)^\top
\phi(\mathbf{x}_j)$. Because $\phi$ is never computed explicitly, $K$ must
also be \textbf{centered in feature space} without access to the (unknown)
feature-space mean:

\[
  \tilde{K} = K - \mathbf{1}_n K - K \mathbf{1}_n + \mathbf{1}_n K \mathbf{1}_n,
\]

where $\mathbf{1}_n$ is the $n \times n$ matrix with every entry $1/n$.
This identity centers $K$ exactly as if $\phi(\mathbf{x}_i)$ had been
mean-subtracted before forming the Gram matrix.

Eigendecomposing $\tilde{K}$ gives eigenvalues $\lambda_1 \ge \dots \ge
\lambda_n \ge 0$ and eigenvectors $\mathbf{a}_1, \dots, \mathbf{a}_n$. The
projection of point $\mathbf{x}$ onto the $i$-th kernel principal component
is then

\[
  z_i(\mathbf{x}) = \sum_{j=1}^{n} a_{ij} \, k(\mathbf{x}_j, \mathbf{x}),
\]

i.e.\ a weighted sum of kernel evaluations against the training points,
normalized so that $\lVert \mathbf{v}_i \rVert = 1$ in feature space
(equivalent to scaling $\mathbf{a}_i$ by $1/\sqrt{\lambda_i}$). Note there
is no explicit $\mathbf{v}_i$ to compute or store --- the component only
ever appears through kernel evaluations, which is exactly what makes the
method tractable in an infinite-dimensional space.

\subsection{Algorithm}

\begin{algorithm}[h]
\caption{Kernel PCA}
\begin{algorithmic}[1]
\State Choose a kernel $k$ (and its hyperparameters, e.g.\ $\gamma$ for RBF)
\State Compute the Gram matrix $K_{ij} = k(\mathbf{x}_i, \mathbf{x}_j)$ for all training pairs
\State Center $K$ in feature space: $\tilde{K} = K - \mathbf{1}_n K - K \mathbf{1}_n + \mathbf{1}_n K \mathbf{1}_n$
\State Compute eigenvalues $\lambda_1 \ge \dots \ge \lambda_n$ and eigenvectors $\mathbf{a}_1, \dots, \mathbf{a}_n$ of $\tilde{K}$
\State Normalize each $\mathbf{a}_i \leftarrow \mathbf{a}_i / \sqrt{\lambda_i}$
\State Choose $k$ components; project new points via $z_i(\mathbf{x}) = \sum_j a_{ij} \, k(\mathbf{x}_j, \mathbf{x})$
\end{algorithmic}
\end{algorithm}

\section{PCA vs. Kernel PCA}

\begin{center}
\begin{tabular}{lll}
  \toprule
  & PCA & Kernel PCA \\
  \midrule
  Structure captured & Linear & Nonlinear (via $\phi$) \\
  Matrix diagonalized & $\Sigma \in \mathbb{R}^{p \times p}$ (covariance) & $\tilde{K} \in \mathbb{R}^{n \times n}$ (centered Gram) \\
  Cost driver & Number of features $p$ & Number of training points $n$ \\
  Components stored as & Vectors $\mathbf{v}_i \in \mathbb{R}^p$ & Weighted combinations of training points \\
  Inverse transform & Exact (linear) & Only approximate (pre-image problem) \\
  \bottomrule
\end{tabular}
\end{center}

Because Kernel PCA diagonalizes an $n \times n$ matrix rather than a
$p \times p$ one, it scales with the number of \emph{samples}, not
features --- the opposite of ordinary PCA, and a relevant consideration for
large training sets.

In the accompanying notebook, Kernel PCA with an RBF kernel is applied to
the same Wine dataset used for PCA/LDA in Part 10 (13 features, 3
cultivar classes), reducing to 2 components before logistic regression.
It reaches the same perfect test accuracy that LDA did, despite being
unsupervised like PCA --- here the nonlinear RBF mapping is expressive
enough to make the classes linearly separable in the transformed space,
something linear PCA could not guarantee.

\section{Key Takeaways}

\begin{itemize}
  \item Kernel PCA applies the kernel trick to PCA: it implicitly maps
        data into a higher-dimensional feature space via $\phi$ and finds
        principal components there, without ever computing $\phi$
        explicitly.
  \item The trick works because PCA can be rewritten entirely in terms of
        pairwise inner products, which the kernel function $k(\mathbf{x}_i,
        \mathbf{x}_j) = \phi(\mathbf{x}_i)^\top \phi(\mathbf{x}_j)$ supplies
        directly.
  \item The RBF kernel is the most common default and corresponds to an
        infinite-dimensional feature space; the choice of kernel and its
        hyperparameters (e.g.\ $\gamma$) controls what kind of nonlinear
        structure can be captured.
  \item Unlike PCA, Kernel PCA's cost scales with the number of training
        points $n$ (via the $n \times n$ Gram matrix), not the number of
        features $p$, and its components have no clean linear form ---
        they only exist as weighted sums of kernel evaluations against the
        training data.
\end{itemize}

\end{document}
